% CS344 Progress Report -- initial draft, 8 November 2026
% Author: Kaya Kaya, Student ID 5587546
%
% Cover layout adapted from University of Warwick Report Template
% by Chris Quinn (2020), as used in specification.tex.
% Source: https://www.overleaf.com/latex/templates/university-of-warwick-report-template/kfpfvcbkdcdk
% Licence: https://creativecommons.org/licenses/by/4.0/
% Changes: progress-report content, subtitle, draft label, margins and headings.
% warwick-crest.pdf is unchanged; University identity and trademark rights
% remain with their respective owners.
%
% Compile with pdfLaTeX twice, from this directory or the repository root.
% Keep assets/warwick-crest.pdf in its subdirectory when uploading to Overleaf.
% No separate class or bibliography file is required.
% This is a starter, not a completed submission. Replace all draft notes.

\documentclass[11pt,a4paper,oneside]{article}

\usepackage[T1]{fontenc}
\usepackage[utf8]{inputenc}
\usepackage[british]{babel}
\usepackage[margin=25mm]{geometry}
\usepackage{amsmath}
\usepackage{amssymb}
\usepackage{graphicx}
\usepackage{booktabs}
\usepackage{tabularx}
\usepackage{array}
\usepackage{microtype}
\usepackage{fancyhdr}
\usepackage[hyphens]{url}
\usepackage[hidelinks,breaklinks]{hyperref}

\graphicspath{{assets/}{docs/reports/assets/}}
\hypersetup{
  pdftitle={Parallel Job Scheduling: Progress Report},
  pdfauthor={Kaya Kaya},
  pdfsubject={CS344 progress-report draft}
}
\pagestyle{fancy}
\fancyhf{}
\fancyhead[L]{CS344: Parallel Job Scheduling}
\fancyhead[R]{Progress report draft}
\fancyfoot[C]{\thepage}
\setlength{\headheight}{14pt}
\setlength{\parindent}{0pt}
\setlength{\parskip}{0.5em}
\setlength{\emergencystretch}{2em}
\urlstyle{same}

% Visible notes identify work still needed before submission.
\newcommand{\draftnote}[1]{\par\textit{To complete: #1}\par}

\begin{document}

\begin{titlepage}
  \centering
  \includegraphics[width=0.35\textwidth]{warwick-crest.pdf}\par
  \vspace{1.3cm}
  {\LARGE\bfseries Parallel Job Scheduling\par}
  \vspace{0.3cm}
  {\large Moldable Jobs in a Loss System\par}
  \vspace{0.9cm}
  {\large Progress Report\par}
  \vspace{0.3cm}
  {\normalsize Initial draft: 8 November 2026\par}
  \vspace{1.2cm}
  {\large\bfseries Kaya Kaya\par}
  \vspace{0.3cm}
  Student ID 5587546\par
  Discrete Mathematics\par
  \vfill
  {\large Supervisor: Dr Arpan Mukhopadhyay\par}
  \vspace{0.4cm}
  {\large\bfseries Department of Computer Science\par}
  \vspace{0.3cm}
  {\large University of Warwick\par}
\end{titlepage}

\setcounter{page}{1}

\section{Introduction and objectives}

Parallel jobs can use several servers at once. Giving more servers to one job
reduces its execution time, but leaves fewer servers available for later arrivals.
This project studies that trade-off in a loss system: a job that finds no idle
server is rejected rather than queued.

The first goal is to fully reproduce the paper by Mukhopadhyay
et al.~\cite{mukhopadhyay2026}: its model, algorithms, mathematical results and
experiments. After completion and review, the supervisor will be asked whether
a finance application is feasible. That later direction remains undecided.

\subsection{Relationship to the original specification}

The original specification proposed order routing between trading venues, with
cancellations and execution deadlines. The current work reproduces the paper's
moldable-job loss system. Financial applicability will be discussed after
reproduction is complete.
The original specification is retained as a record of the initial proposal.

\draftnote{Record the reasons for the reproduction-first sequence, its effect on
the original plan, and supervisor feedback. Finance applicability remains undecided.}

\section{Background and model}

A \emph{moldable job} can use a chosen number of servers, fixed from its start
until completion. If its one-server execution time is $X$ and its allocation is
$k$, its execution time is
\begin{equation}
  T = \frac{X}{s_k},
\end{equation}
where $s_k$ is the speedup on $k$ servers. A concave speedup curve represents
diminishing gains from additional servers.

The paper derives an allocation optimisation problem and shows that an optimal
solution can use at most two neighbouring allocation sizes under its assumptions.
Its allocation policy has asymptotic optimality guarantees as the system grows;
this does not establish that it is the best policy for every finite system
\cite{mukhopadhyay2026}.

\draftnote{Explain the optimisation, its assumptions and benchmark. Compare the
paper critically with related work, including differences in queueing and speedup
assumptions.}

\subsection{Current baseline}

The initial implementation uses the following single-class model.

\begin{center}
\begin{tabularx}{\linewidth}{@{}lX@{}}
\toprule
Setting & Baseline value \\
\midrule
Servers & $n=100$, identical and initially idle \\
Maximum allocation & $d=5$ servers per job \\
Job sizes & Independent exponential $X$, with $\mathbb{E}[X]=1$ \\
Arrivals & Poisson, independent of job sizes \\
Normalized load & $\rho=0.8$ \\
Arrival rate & $a=\rho n/\mathbb{E}[X]=80$ per time unit \\
Speedups & $(s_1,s_2,s_3,s_4,s_5)=(1,1.8,2.5,3,3.4)$ \\
Waiting queue & None \\
\bottomrule
\end{tabularx}
\end{center}

Greedy allocation assigns $\min(d,i)$ servers when $i>0$ servers are idle.
For this baseline, the randomized policy samples a target of three servers with probability
$0.625$, or four with probability $0.375$, and assigns the smaller of that target
and $i$. Both policies reject an arrival only when $i=0$.

For each randomized experiment, separate helpers calculate the optimal target
counts and their probabilities from the configured load and speedup curve.
The main entry point passes this distribution to the allocation policy, which
samples a target and caps it by idle capacity.

\section{Implementation and validation}

An event-driven Python simulator advances between arrivals and completions.
Completion events are kept in a priority queue. This avoids stepping through time
intervals in which nothing happens. Completions are processed before arrivals at
the same time, and all servers allocated to a completed job are released together.

The implementation separates immutable job inputs, server capacity, simulation control,
allocation policies, workload generation and statistics. Reusable code is in the
\texttt{scheduling/} package. A single \texttt{main.py} entry point runs one seeded
simulation of the randomized policy. Allocation decisions use
a separate random stream.

Automated checks cover deterministic event traces, server conservation, rejection
when capacity is exhausted, simultaneous events, measurement cutoffs and seeded
repeatability. These checks support
implementation correctness; they do not establish agreement with all of the
paper's experiments.

\draftnote{Add a worked validation example and explain the main design choices
and limitations using specific evidence.}

\section{Evaluation method and current evidence}

For measured arrivals, the blocking estimate is
\begin{equation}
  \widehat{P}_{\mathrm{block}}
  = \frac{N_{\mathrm{rejected}}}{N_{\mathrm{arrivals}}}.
\end{equation}
The mean execution-time estimate includes accepted jobs only:
\begin{equation}
  \widehat{D}
  = \frac{1}{N_{\mathrm{accepted}}}
    \sum_{j\,\mathrm{accepted}} T_j.
\end{equation}
The first estimate is undefined if there are no measured arrivals; the second is
undefined if none is accepted.

A configurable initial number of arrivals is excluded as warm-up. Active jobs
remain in the system at the cutoff. Measured accepted jobs contribute their full
execution times, including completions after the last arrival.

The current entry point runs one seeded randomized-policy simulation. Earlier pilot output
files were removed during repository cleanup on 8 November; no saved experiment
results are currently available. Numerical agreement with Figure 2 and Table I
has not yet been established.

\draftnote{Specify the repeated-run protocol, justify warm-up and run lengths,
report uncertainty, and add saved results with settings and reproduction commands.
Resolve any unreported paper settings before claiming numerical reproduction.}

\section{Remaining work and timetable}

The next tasks are to derive the arrival rate from model settings, calculate the
optimal allocation targets and their probabilities, and validate those
calculations against known cases. Subsequent work will add repeated experiments,
reproduce Figure 2 and Table I, and implement the heterogeneous-workload model
in Section VI of the paper.

The progress-report deadline is 10 December 2026 at 12 noon. The available budget
until then is approximately eight to ten hours per week. Mathematical reading,
report writing and supervisor feedback must run alongside implementation.

\draftnote{Add a feasible updated timetable with dependencies, effort estimates,
contingency and an agreed supervisor-feedback date. Distinguish work planned
before the progress report from work planned for the rest of the project.}

\section{Reflection and professional considerations}

The current implementation provides a baseline, while general allocation
calculations, complete experimental evidence and the multi-class extension remain
unfinished. Reproduction also depends on clarifying missing experimental details,
including the Mixed-Erlang construction and the original measurement conventions.
These questions should be resolved with the supervisor or recorded as explicit
assumptions.

\draftnote{Reflect on specific difficulties, decisions and changes to the plan.
Explain relevant legal, social, ethical and professional issues. Record the
ethical-review status after checking the applicable requirements.}

\begin{thebibliography}{9}
\bibitem{mukhopadhyay2026}
A. Mukhopadhyay, S. Ghanbarian, R. R. Mazumdar and F. Guillemin,
``On Optimal Server Allocation for a Loss System With Moldable Jobs,''
\emph{IEEE Transactions on Networking}, vol.~34, pp.~384--397, 2026.
doi: \href{https://doi.org/10.1109/TON.2025.3598595}{10.1109/TON.2025.3598595}.
\end{thebibliography}

\appendix
\section{Original specification}

\draftnote{Append the original submitted specification to the final submission
PDF. Include an approved revised specification separately if appropriate.}

\end{document}
